% ============================================================
%  补充宏包与自定义命令
%  注意：页面、表格、代码、颜色等核心配置见 preamble.tex
% ============================================================

% --- 数学公式 ---
\usepackage{amsmath}

% --- 浮动体强制放置 [H] ---
\usepackage{float}

% --- 边距调整（\begin{addmargin}...\end{addmargin}）---
\usepackage{scrextend}

% --- 多栏排版 ---
\usepackage{multicol}

% --- 按名称引用（\nameref）---
\usepackage{nameref}

% --- 附录 ---
\usepackage{appendix}

% --- TODO 便签（撰写阶段使用，定稿后可注释）---
\usepackage{xargs}
\usepackage[colorinlistoftodos,prependcaption,textsize=footnotesize]{todonotes}
\newcommandx{\commred}[2][1=]{\textcolor{Red}
{\todo[linecolor=red,backgroundcolor=red!25,bordercolor=red,#1]{#2}}}
\newcommandx{\commblue}[2][1=]{\textcolor{Blue}
{\todo[linecolor=blue,backgroundcolor=blue!25,bordercolor=blue,#1]{#2}}}
\newcommandx{\commgreen}[2][1=]{\textcolor{OliveGreen}
{\todo[linecolor=OliveGreen,backgroundcolor=OliveGreen!25,bordercolor=OliveGreen,#1]{#2}}}
\newcommandx{\commpurp}[2][1=]{\textcolor{Plum}
{\todo[linecolor=Plum,backgroundcolor=Plum!25,bordercolor=Plum,#1]{#2}}}

% --- 简单间距 ---
\newcommand\tab[1][1cm]{\hspace*{#1}}

% --- 等宽字体快捷方式 ---
\def\code#1{{\tt #1}}
\def\note#1{\noindent{\bf [Note: #1]}}

% --- 附录编号格式（附录 A、附录 B…）---
\makeatletter
\def\@seccntformat#1{\@ifundefined{#1@cntformat}%
   {\csname the#1\endcsname\quad}%  默认
   {\csname #1@cntformat\endcsname}% 单独控制
}
\let\oldappendix\appendix
\renewcommand\appendix{%
    \oldappendix
    \newcommand{\section@cntformat}{\appendixname~\thesection\quad}
}
\makeatother
